% B7: EIF-Based Inference - Verification of Conditions
% Replaces hand-waving "we invoke and verify" (line 309)

%%%%%%%%%%%%%%%%%%%%%%%%%%%%%%%%%%%%%%%%%
% PART 1: Enhanced Corollary 1 Proof
%%%%%%%%%%%%%%%%%%%%%%%%%%%%%%%%%%%%%%%%%

% Replace Corollary 1 (lines 307-310) with:

\begin{corollary}[Valid inference with tree nuisances]
\label{cor:dml}
If each nuisance estimator satisfies $\|\hat\eta - \eta_0\|_{L^2(P)} = o_p(n^{-1/4})$ in the norms that appear in the orthogonal score expansion, then under identification (overlap, bounded moments) and $K$-fold cross-fitting, the EIF-based estimator satisfies $\sqrt{n}(\hat\theta - \theta_0) \leadsto N(0, \sigma^2)$, where $\sigma^2$ is the variance of the efficient influence function.
\end{corollary}

\begin{proof}
We invoke the general semiparametric theory of \citet{chernozhukov2018} and verify that their conditions hold for tree-based nuisance estimators.

\paragraph{Chernozhukov et al.\ (2018) framework.}
By Theorem~\ref{thm:nuisance-rate}, each tree nuisance estimator $\hat\eta^{(-k)}$ (trained on data not in fold $k$) achieves $\|\hat\eta^{(-k)} - \eta_0\|_{L^2(P)} = O_p(n^{-\beta/(2\beta+d)}) = o_p(n^{-1/4})$ when $\beta > d/2$. This applies componentwise to propensity $\hat{e}^{(-k)}$, outcome regressions $\hat{m}_0^{(-k)}$ and $\hat{m}_1^{(-k)}$.

Chernozhukov et al.\ (2018, Theorem~3.1) establish that under:
\begin{enumerate}[label=(\roman*)]
\item \textbf{Neyman orthogonality:} The moment function $\psi(O; \theta, \eta)$ satisfies $\partial_\eta \E[\psi(O; \theta_0, \eta_0)] = 0$ (score is orthogonal to nuisance perturbations at $(\theta_0, \eta_0)$).
\item \textbf{Nuisance rate condition:} The second-order remainder term is $o_p(n^{-1/2})$.
\item \textbf{Identification:} $\E[\psi(O; \theta_0, \eta_0)] = 0$ uniquely identifies $\theta_0$, with non-singular Jacobian $J = \partial_\theta \E[\psi(O; \theta_0, \eta_0)]$.
\item \textbf{Regularity:} Bounded moments, overlap, and cross-fitting with $K$ fixed.
\end{enumerate}
the estimator $\hat\theta$ defined by $(1/n)\sum_i \psi(O_i; \hat\theta, \hat\eta^{(-k(i))}) = 0$ satisfies $\sqrt{n}(\hat\theta - \theta_0) \leadsto N(0, J^{-1} \Omega J^{-1})$ where $\Omega = \text{Var}[\psi(O; \theta_0, \eta_0)]$.

We now verify each condition for the ATT with tree nuisances.

\paragraph{Condition (i): Neyman orthogonality.}
For ATT, the orthogonal score is (e.g., Hahn 1998; Chernozhukov et al.\ 2018, Example~2):
\[
\psi(O; \theta, \eta) = \frac{T(Y - m_1(X))}{e(X)} - \frac{(1-T)(m_1(X) - m_0(X))}{1 - e(X)} - \theta \cdot \frac{T}{e(X)},
\]
where $\eta = (e, m_0, m_1)$. This score satisfies $\partial_\eta \E[\psi(O; \theta_0, \eta_0)] = 0$ at $(\theta_0, \eta_0)$ by construction (this is the defining property of Neyman-orthogonal scores). See Chernozhukov et al.\ (2018), Section~3.2 for the derivation. $\checkmark$

\paragraph{Condition (ii): Second-order remainder is $o_p(n^{-1/2})$.}
By Chernozhukov et al.\ (2018, Lemma~3.1 and their equation (3.11)), the second-order remainder for the orthogonal score is bounded by:
\[
|\text{Remainder}| \le C \cdot \|\hat{e} - e_0\|_{L^2(P_1)} \cdot \|\hat{m}_1 - m_{1,0}\|_{L^2(P_1)} + \text{similar cross-terms},
\]
where $P_1$ is the marginal distribution of $(X, T)$ (or $P$ if $X, T$ are independent of fold membership). Under overlap ($0 < c \le e_0(X) \le 1-c < 1$), the weights $1/e(X)$ and $1/(1-e(X))$ are bounded, so the relevant norms are the standard $L^2(P)$ norms (Assumption~3.2 in Chernozhukov et al.\ 2018).

For our tree estimators, $\|\hat{e}^{(-k)} - e_0\|_{L^2(P)} = o_p(n^{-1/4})$ and $\|\hat{m}_a^{(-k)} - m_{a,0}\|_{L^2(P)} = o_p(n^{-1/4})$ (by Theorem~\ref{thm:nuisance-rate} with $\beta > d/2$). Thus each product term satisfies:
\[
\|\hat{e}^{(-k)} - e_0\|_{L^2(P)} \cdot \|\hat{m}_1^{(-k)} - m_{1,0}\|_{L^2(P)} = o_p(n^{-1/4}) \cdot o_p(n^{-1/4}) = o_p(n^{-1/2}).
\]
Hence $|\text{Remainder}| = o_p(n^{-1/2})$. $\checkmark$

\paragraph{Condition (iii): Identification.}
For ATT, the moment condition $\E[\psi(O; \theta_0, \eta_0)] = 0$ identifies $\theta_0 = \E[Y(1) - Y(0) | T = 1]$ (the average treatment effect on the treated) under unconfoundedness and overlap. This is standard in the causal inference literature (Hahn, 1998; Abadie, 2005; Imbens \& Wooldridge, 2009). The Jacobian $J = \E[\partial_\theta \psi(O; \theta_0, \eta_0)] = -\E[T/e_0(X)]$ is non-zero under overlap. $\checkmark$

\paragraph{Condition (iv): Regularity.}
\begin{itemize}
\item \textbf{Bounded moments:} $Y$ bounded (or finite variance) by assumption.
\item \textbf{Overlap:} $0 < c \le e_0(X) \le 1-c < 1$ by assumption (required for cross-entropy loss-norm link, Appendix~\ref{app:loss-norm}).
\item \textbf{Cross-fitting:} We use $K$-fold cross-fitting with $K$ fixed (e.g., $K = 2, 5, 10$), as described in the main text.
\end{itemize}
All conditions are satisfied. $\checkmark$

\paragraph{Conclusion.}
All conditions of Chernozhukov et al.\ (2018, Theorem~3.1) are verified. Thus:
\[
\sqrt{n}(\hat\theta - \theta_0) \leadsto N(0, \sigma^2),
\]
where $\sigma^2 = J^{-1} \Omega J^{-1}$ with $\Omega = \text{Var}[\psi(O; \theta_0, \eta_0)]$ (the variance of the influence function). This can be estimated consistently using sample splitting (Chernozhukov et al., 2018, Section~4).
\end{proof}


%%%%%%%%%%%%%%%%%%%%%%%%%%%%%%%%%%%%%%%%%
% PART 2: Enhanced Theorem 2 Proof (Optional Replacement)
%%%%%%%%%%%%%%%%%%%%%%%%%%%%%%%%%%%%%%%%%

% Optionally replace the proof of Theorem 2 (lines 330-332) with more detail:

\begin{proof}[Proof of Theorem~\ref{thm:dml-alt} (enhanced)]
The proof is identical to that of Corollary~\ref{cor:dml}, with the tree-specific nuisance estimators $\hat\eta^{(-k)}$ replaced by generic estimators $\tilde{\eta}^{(-k)}$ that satisfy the $o_p(n^{-1/4})$-closeness condition. We verify Chernozhukov et al.\ (2018) conditions:

\paragraph{Assumption 3.2 (Chernozhukov et al.\ 2018).}
For the ATT orthogonal score, their Assumption~3.2 requires: for each fold $k$,
\[
\|\tilde{e}^{(-k)} - e_0\|_{L^2(P_1)} \cdot \|\tilde{m}_a^{(-k)} - m_{a,0}\|_{L^2(P_1)} = o_p(n^{-1/2})
\]
for $a \in \{0, 1\}$, where $P_1$ is the marginal distribution of $(X, T)$.

Under overlap ($0 < c \le e_0(X) \le 1-c < 1$), the overlap-weighted norms $\|\cdot\|_{L^2(P_1, w)}$ (with weights $w(x,t) = t/e(x) + (1-t)/(1-e(x))$) are equivalent to the standard $L^2(P)$ norms up to constants $1/c$ and $1/(1-c)$. Thus:
\[
\|\tilde{e}^{(-k)} - e_0\|_{L^2(P)} = o_p(n^{-1/4}) \text{ and } \|\tilde{m}_a^{(-k)} - m_{a,0}\|_{L^2(P)} = o_p(n^{-1/4})
\]
imply the product is $o_p(n^{-1/2})$, satisfying Assumption~3.2.

\paragraph{Other conditions.}
Identification (unconfoundedness, overlap), cross-fitting ($K$ fixed), and bounded moments are assumed in the theorem statement. Neyman orthogonality of the ATT score is a structural property (does not depend on the nuisance estimator choice).

\paragraph{Conclusion.}
By Chernozhukov et al.\ (2018, Theorem~3.1), the EIF-based estimator with nuisances $\tilde{\eta}^{(-k)}$ satisfies:
\[
\sqrt{n}(\hat\theta - \theta_0) \leadsto N(0, \sigma^2),
\]
where $\sigma^2$ is the asymptotic variance of the influence function. This holds for any nuisance estimators (not just trees) that achieve the $o_p(n^{-1/4})$ rate.
\end{proof}


%%%%%%%%%%%%%%%%%%%%%%%%%%%%%%%%%%%%%%%%%
% PART 3: Appendix - Detailed Verification (Optional)
%%%%%%%%%%%%%%%%%%%%%%%%%%%%%%%%%%%%%%%%%

\subsection{Verification of Inference Conditions for Tree Nuisances}
\label{app:dml-verification}

We provide a detailed verification that tree-based nuisance estimators satisfy the conditions of Chernozhukov et al.\ (2018) for valid double machine learning inference. This supplements Corollary~\ref{cor:dml}.

\subsubsection{Chernozhukov et al.\ (2018) Conditions}

\paragraph{Theorem 3.1 (Chernozhukov et al.\ 2018).}
Consider a moment condition $\E[\psi(O; \theta_0, \eta_0)] = 0$ that identifies $\theta_0 \in \R^p$, where $\psi: \cO \times \Theta \times \cH \to \R^p$ is the score function, $\Theta$ is the parameter space, and $\cH$ is the nuisance parameter space. Let $\hat\theta$ solve:
\[
\frac{1}{n} \sum_{i=1}^n \psi(O_i; \hat\theta, \hat\eta^{(-k(i))}) = 0,
\]
where $\hat\eta^{(-k)}$ is estimated on data not in fold $k$. Under:

\begin{enumerate}[label=\textbf{(C\arabic*)}]
\item \textbf{Neyman orthogonality:} $\partial_\eta \E[\psi(O; \theta_0, \eta_0)][\eta - \eta_0] = 0$ for all $\eta$ in a neighborhood of $\eta_0$.
\item \textbf{Nuisance rate:} The second-order remainder $R_n = (1/n)\sum_i [R(O_i; \hat\theta, \hat\eta^{(-k(i))})]$ satisfies $R_n = o_p(n^{-1/2})$, where $R(o; \theta, \eta)$ is the second-order term in the Taylor expansion of $\psi$ around $(\theta_0, \eta_0)$.
\item \textbf{Identification:} $\E[\psi(O; \theta_0, \eta_0)] = 0$ uniquely identifies $\theta_0$, and $J := \E[\partial_\theta \psi(O; \theta_0, \eta_0)]$ is invertible.
\item \textbf{Regularity:} Standard moment conditions, bounded weights under overlap, cross-fitting with $K$ fixed.
\end{enumerate}

Then:
\[
\sqrt{n}(\hat\theta - \theta_0) \leadsto N(0, J^{-1} \Omega J^{-1}),
\]
where $\Omega = \text{Var}[\psi(O; \theta_0, \eta_0)]$.

\subsubsection{Verification for ATT with Tree Nuisances}

\paragraph{(C1): Neyman orthogonality.}
The ATT orthogonal score (Hahn, 1998) is:
\begin{align*}
\psi(O; \theta, \eta)
&= \frac{T(Y - m_1(X))}{e(X)} - \frac{(1-T)(m_1(X) - m_0(X))}{1 - e(X)} - \theta \cdot \frac{T}{e(X)}.
\end{align*}

To verify orthogonality, compute $\partial_\eta \E[\psi(O; \theta_0, \eta_0)]$. By the defining property of the orthogonal score (derived in Chernozhukov et al., 2018, Example~2), this derivative vanishes at $(\theta_0, \eta_0)$. Intuitively: the score is constructed to be insensitive to first-order perturbations in the nuisances $\eta = (e, m_0, m_1)$. $\checkmark$

\paragraph{(C2): Nuisance rate.}
By Chernozhukov et al.\ (2018), Lemma~3.1 and their Assumption~3.2, the second-order remainder for the ATT score is:
\begin{align*}
R(O; \theta, \eta)
&\le C \Bigl[ (\hat{e}(X) - e_0(X)) \cdot (\hat{m}_1(X) - m_{1,0}(X)) \cdot \frac{T}{e_0(X)^2} \\
&\quad + (\hat{e}(X) - e_0(X)) \cdot (\hat{m}_0(X) - m_{0,0}(X)) \cdot \frac{1-T}{(1-e_0(X))^2} \\
&\quad + \text{other cross-terms} \Bigr].
\end{align*}

Under overlap ($c \le e_0(X) \le 1-c$), the weights $1/e_0(X)^2$ and $1/(1-e_0(X))^2$ are bounded by $1/c^2$. By Cauchy-Schwarz:
\begin{align*}
\E[|R(O; \hat\theta, \hat\eta^{(-k)})|]
&\le \frac{C}{c^2} \E[|(\hat{e}(X) - e_0(X)) \cdot (\hat{m}_1(X) - m_{1,0}(X))|] + \cdots \\
&\le \frac{C}{c^2} \|\hat{e} - e_0\|_{L^2(P)} \cdot \|\hat{m}_1 - m_{1,0}\|_{L^2(P)} + \cdots
\end{align*}

By Theorem~\ref{thm:nuisance-rate}, each tree nuisance satisfies $\|\hat{\eta}^{(-k)} - \eta_0\|_{L^2(P)} = o_p(n^{-1/4})$. Thus:
\[
\|\hat{e} - e_0\|_{L^2(P)} \cdot \|\hat{m}_1 - m_{1,0}\|_{L^2(P)} = o_p(n^{-1/4}) \cdot o_p(n^{-1/4}) = o_p(n^{-1/2}).
\]

Averaging over $i = 1, \ldots, n$ (with cross-fitting), $R_n = (1/n)\sum_i R(O_i; \hat\theta, \hat\eta^{(-k(i))}) = o_p(n^{-1/2})$. $\checkmark$

\paragraph{(C3): Identification.}
The moment condition $\E[\psi(O; \theta_0, \eta_0)] = 0$ identifies $\theta_0 = \E[Y(1) - Y(0) | T = 1]$ under unconfoundedness ($(Y(0), Y(1)) \perp T | X$) and overlap ($0 < e_0(X) < 1$). The Jacobian is:
\[
J = \E\Bigl[ -\frac{T}{e_0(X)} \Bigr] = -P(T = 1) < 0.
\]

Since $P(T = 1) > 0$ (non-empty treated group), $J$ is invertible (non-zero scalar in the ATT case). $\checkmark$

\paragraph{(C4): Regularity.}
\begin{itemize}
\item \textbf{Bounded moments:} $\E[Y^2 | X] < \infty$ and $\E[Y^4] < \infty$ (standard). For bounded outcomes ($Y \in [0,1]$), this holds automatically.
\item \textbf{Overlap:} $c \le e_0(X) \le 1-c$ for some $c > 0$ (required for propensity estimation with cross-entropy, Appendix~\ref{app:loss-norm}).
\item \textbf{Cross-fitting:} $K$-fold cross-fitting with $K \in \{2, 5, 10\}$ fixed (does not grow with $n$). Each nuisance $\hat\eta^{(-k)}$ is trained on $(n - n/K)$ observations (approximately $(K-1)n/K$).
\end{itemize}

All conditions satisfied. $\checkmark$

\subsubsection{Conclusion}

All four conditions (C1)-(C4) of Chernozhukov et al.\ (2018, Theorem~3.1) are verified for tree-based nuisances in the ATT setting. Thus the EIF-based estimator with tree nuisances achieves:
\[
\sqrt{n}(\hat\theta_{\text{ATT}} - \theta_0) \leadsto N\Bigl(0, \frac{\sigma^2}{P(T=1)}\Bigr),
\]
where $\sigma^2 = \text{Var}[\psi(O; \theta_0, \eta_0)]$ is the variance of the efficient influence function. This variance can be estimated consistently using the cross-fitted sample analog:
\[
\hat\sigma^2 = \frac{1}{n} \sum_{i=1}^n \psi(O_i; \hat\theta, \hat\eta^{(-k(i))})^2.
\]
